% from main.tex
%%%%%%%%%%  scan idx 122  |  folio 109  %%%%%%%%%%
suggérer qu'il s'agit de \(A\)-Modules à gauche), ou même \(\Dpx(S)\) s'il n'y a pas d'ambiguïté sur \(A\). On notera \(S\) l'objet final de \(S\) et posera \(\Dpx(\mathcal C_S)=\Dpx({}_A S)\) (ou \(\Dpx(S)\)). Un objet de \(\Dpx(S)\) sera dit \(n\)-pseudo-cohérent (resp. pseudo-cohérent) s'il est \(n\)-\(\mathcal C_0\)-pseudo-cohérent (resp. \(\mathcal C_0\)-pseudo-cohérent) (itou pour les notions strictes). La catégorie fibrée \(\Dpx(\mathcal C)_{n-\mathcal C_0-\coh}\) (resp. \(\Dpx(\mathcal C)_{\mathcal C_0-\coh}\)) (2.6) sera notée \(\Dpx({}_A S)_{n-\coh}\) (resp. \(\Dpx({}_A S)_{\coh}\)). On observera qu'en vertu de (2.14) on obtient la même notion de pseudo-cohérence si l'on remplace \(\mathcal C_0\) par la \(S\)-catégorie \(\mathcal C_1\) définie par \(U\longmapsto\) (catégorie des \(A_U\)-Modules à gauche localement facteurs directs de Modules libres de type fini). Si un \(A\)-Module \(M\) est facteur direct d'un \(A\)-Module libre de type fini, il ne s'ensuit pas, en général, que \(M\) est localement libre de type fini. On a cependant :

\textbf{Lemme 2.15.1.} \emph{Supposons que \(S\) vérifie l'une des deux hypothèses suivantes :}

\emph{(i) \(S\) « a assez de points » (SGA 4 IV 6.4) et pour tout point \(s\) de \(S\) les \(A_s\)-modules (à gauche) projectifs de type fini sont libres, cette dernière condition étant réalisée par exemple (BOURBAKI Alg. Comm. chap. II \S~3 n\textsuperscript{o}~2 cor.~2 de la prop. 5) si \(A_s/\operatorname{radical}(A_s)\) est un corps ou si \(\operatorname{radical}(A_s)\) est nilpotent.}

\emph{(ii) \((S,A)\) est un topos « localement annelé » au sens de [1], ce qui signifie que \(A\) est commutatif et que pour tout \(X\in\Ob S\) et tout \(f\in\Gamma(X,A)\) on a \(X=X_f\cup X_{1-f}\), où \(X_f\) (resp. \(X_{1-f}\)) désigne le plus grand sous-objet de \(X\) où \(f\) (resp. \(1-f\)) est inversible.}

\emph{Alors tout facteur direct d'un \(A\)-module (à gauche) libre de type fini est localement libre de type fini.}

\emph{Preuve.} Supposons (i). On remarque d'abord que si \(E\) et \(F\) sont deux \(A\)-Modules, \(E\) étant de présentation finie, la flèche canonique
